\section{Introduction}
3D scene reconstruction and novel view synthesis (NVS) are rapidly advancing, with many applications across different domains~\cite{EVALScenReconstruction2024}.
These methods can be applied in fields such as Virtual Reality (VR) for creating immersive worlds, cinematography to create visually appealing assets efficiently, or
robot vision to help robots understand their physical environment~\cite{SLAMROBOT2021}.
The foundation of 3D scene reconstruction was laid by traditional Structure from Motion pipelines. More recently, significant advances in NVS were achieved by representing the scene as Neural Radiance Fields (NeRF)~\cite{NeRFOriginal}. These methods achieve state-of-the-art results but suffer from slow training speeds and are unsuitable for real-time rendering due to high latency. Newer methods such as 3D Gaussian Splatting (3DGS)~\cite{3dgs2023} enable high-quality NVS even for real-time rendering. Additionally, training speed is significantly reduced.

A major limitation of these novel view synthesis methods is the need for dense views, which often is not feasible for real-world applications. In sparse-view settings, these methods tend to struggle with bad initialization, depth ambiguities and overfitting to training views. To improve the performance in these settings and counteract the depth ambiguities, certain priors are introduced to better generalize and escape minima throughout the optimization process. Methods such as \textit{DNGaussian}~\cite{DNGaussian} and \textit{Few-shot Novel View Synthesis using Depth}~\cite{few-shotNVS3DGS} leverage monocular depth estimators to regularize the model with the help of the inferred depth. The depth regularization helps significantly to improve the depth ambiguities and increase the generalization capability of the models. A challenge for these models is correct depth scaling, proper point initialization, and accurate camera poses. The initial points and camera poses are traditionally generated using Structure from Motion (SfM) pipelines. However, these pipelines often struggle with sparse input views and limited overlap between views. Recent advancements for sparse-view settings were achieved by leveraging dense initialization with the help of point cloud estimation networks~\cite{InstantSplat, InternGS, SparseGS}. They replace the traditional SfM pipeline with models such as \textit{MASt3R}~\cite{mastrsfm2025} or \textit{DUSt3R}~\cite{dust3r_cvpr24} for the point cloud estimation and camera pose estimation.
The resulting models achieve high-fidelity results but require good initial reference views for accurate predictions. In addition, a time-consuming iterative camera alignment process is required, which can take several minutes. Inaccurate camera poses may further reduce reconstruction quality.

We make the following contributions:
\begin{itemize}
    \item We discuss a method for leveraging a Permutation-Equivariant point cloud estimation network for dense initialization without relying on traditional SfM.
    \item We introduce confidence aware pearson depth loss, to counteract uncertain depth estimations.
    \item We explore the use of PGSR in sparse-view settings for improved geometry alignment and reduced overfitting.
\end{itemize}
Our method achieves state-of-the-art results in sparse-view settings and significantly improves Gaussian surface alignment, while reducing floaters.
Our code is publicly available at \href{https://github.com/Mango0000/Pi-GS}{https://github.com/Mango0000/Pi-GS}.